\documentclass[12pt,a4paper]{exam}
\usepackage[utf8]{inputenc}
\usepackage{amsmath,amssymb,amsfonts}
\usepackage{geometry}
\usepackage{enumitem}
\usepackage{graphicx}
\usepackage{hyperref}
\usepackage{xcolor}
\geometry{a4paper, top=2cm, bottom=2cm, left=2.5cm, right=2.5cm}
\hypersetup{colorlinks=true, linkcolor=blue, urlcolor=blue}
\renewcommand{\thequestion}{\textbf{\arabic{question}}}
\renewcommand{\questionlabel}{\thequestion.}
\pointname{ marks}
\pointsdroppedatright
\bracketedpoints
\setlength{\parindent}{0pt}
\setlength{\emergencystretch}{3em}
\allowdisplaybreaks
\newsavebox{\schmathbox}
\newcommand{\fitmath}[1]{%
  \sbox{\schmathbox}{$\displaystyle #1$}%
  \ifdim\wd\schmathbox>\linewidth
    \resizebox{\linewidth}{!}{\usebox{\schmathbox}}%
  \else
    \usebox{\schmathbox}%
  \fi}

\begin{document}

\thispagestyle{empty}
\begin{center}
  \textbf{\LARGE Diag} \\[0.3cm]
  \textbf{Time Allowed: 3 Hours} \hfill \textbf{Maximum Marks: 80} \\[0.3cm]
  \rule{\textwidth}{0.5pt}
\end{center}

\vspace{0.3cm}

\noindent\textbf{General Instructions:}
\begin{enumerate}
  \item {\normalfont\normalcolor This question paper contains 40 questions. All questions are compulsory.}
  \item {\normalfont\normalcolor Questions are grouped into sections A through E.}
  \item {\normalfont\normalcolor Use of calculators is NOT permitted.}
\end{enumerate}
\bigskip
\noindent\textbf{\large Section A: Multiple Choice \& Assertion-Reason (1 mark each)}

\begin{questions}
\question[1] Assertion (A): The common difference of the Arithmetic Progression $5, 2, -1, -4, \dots$ is $-3$. Reason (R): The common difference $d$ of an AP is given by $a_{n} - a_{n-1}$.
\begin{enumerate}[label=(\Alph*)]
  \item (A) Both A and R are true and R is the correct explanation of A
  \item (B) Both A and R are true but R is NOT the correct explanation of A
  \item (C) A is true but R is false
  \item (D) A is false but R is true
\end{enumerate}
\vspace{0.15cm}

\question[1] Assertion (A): The sum of the first $n$ terms of an arithmetic progression is given by $S_n = \frac{n}{2}[2a + (n-1)d]$, where $a$ is the first term and $d$ is the common difference.
Reason (R): The sum of the first $n$ terms of an AP can be calculated as $S_n = \frac{n}{2}(\text{first term} + \text{last term})$.
\begin{enumerate}[label=(\Alph*)]
  \item (A) Both A and R are true and R is the correct explanation of A
  \item (B) Both A and R are true but R is NOT the correct explanation of A
  \item (C) A is true but R is false
  \item (D) A is false but R is true
\end{enumerate}
\vspace{0.15cm}

\question[1] If $\sin \theta = \frac{3}{5}$, then the value of $\cos \theta$ is:
\begin{enumerate}[label=(\Alph*)]
  \item (A) 2/5
  \item (B) 3/4
  \item (C) 4/5
  \item (D) 5/3
\end{enumerate}
\vspace{0.15cm}

\question[1] The value of $\tan 30^\circ \cdot \cot 30^\circ$ is:
\begin{enumerate}[label=(\Alph*)]
  \item (A) 0
  \item (B) 1
  \item (C) 1/2
  \item $(D) \sqrt{3}$
\end{enumerate}
\vspace{0.15cm}

\question[1] If $A = 60^\circ$ and $B = 30^\circ$, then the value of $\sin(A+B)$ is:
\begin{enumerate}[label=(\Alph*)]
  \item (A) 0
  \item (B) 1/2
  \item $(C) \frac{\sqrt{3}}{2}$
  \item (D) 1
\end{enumerate}
\vspace{0.15cm}

\question[1] A tower stands vertically on the ground. From a point on the ground, which is $15$ m away from the foot of the tower, the angle of elevation of the top of the tower is found to be $60^\circ$. The height of the tower is:
\begin{enumerate}[label=(\Alph*)]
  \item (A) $5\sqrt{3}$ m
  \item (B) $15$ m
  \item (C) $15\sqrt{3}$ m
  \item (D) $30$ m
\end{enumerate}
\vspace{0.15cm}

\question[1] The angle of elevation of the top of a pole from a point on the ground is $30^\circ$. If the pole is $10$ m high, the distance of the point from the foot of the pole is:
\begin{enumerate}[label=(\Alph*)]
  \item (A) $10$ m
  \item (B) $10\sqrt{3}$ m
  \item (C) $\frac{10}{\sqrt{3}}$ m
  \item (D) $20$ m
\end{enumerate}
\vspace{0.15cm}

\question[1] If the length of a string of a kite flying at a height of $60$ m from the ground is $120$ m, then the angle of inclination of the string with the ground is:
\begin{enumerate}[label=(\Alph*)]
  \item (A) $30^\circ$
  \item (B) $45^\circ$
  \item (C) $60^\circ$
  \item (D) $90^\circ$
\end{enumerate}
\vspace{0.15cm}

\question[1] Assertion (A): The common difference of the AP $5, 2, -1, -4, \dots$ is $-3$. Reason (R): The common difference $d$ of an AP $a_1, a_2, a_3, \dots$ is given by $d = a_{n} - a_{n-1}$.
\begin{enumerate}[label=(\Alph*)]
  \item (A) Both Assertion (A) and Reason (R) are true and Reason (R) is the correct explanation of Assertion (A)
  \item (B) Both Assertion (A) and Reason (R) are true but Reason (R) is NOT the correct explanation of Assertion (A)
  \item (C) Assertion (A) is true but Reason (R) is false
  \item (D) Assertion (A) is false but Reason (R) is true
\end{enumerate}
\vspace{0.15cm}

\question[1] Assertion (A): The common difference of an A.P. where $a_n = 3n + 7$ is 3. Reason (R): The common difference of an A.P. given by $a_n = An + B$ is always $A$.
\begin{enumerate}[label=(\Alph*)]
  \item (A) Both Assertion (A) and Reason (R) are true and Reason (R) is the correct explanation of Assertion (A)
  \item (B) Both Assertion (A) and Reason (R) are true but Reason (R) is NOT the correct explanation of Assertion (A)
  \item (C) Assertion (A) is true but Reason (R) is false
  \item (D) Assertion (A) is false but Reason (R) is true
\end{enumerate}
\vspace{0.15cm}

\question[1] Assertion (A): The common difference of the Arithmetic Progression 5, 2, -1, -4, ... is -3. Reason (R): The common difference 'd' of an AP is given by the formula $d = a_{n} - a_{n-1}$.
\begin{enumerate}[label=(\Alph*)]
  \item (A) Both Assertion (A) and Reason (R) are true and Reason (R) is the correct explanation of Assertion (A)
  \item (B) Both Assertion (A) and Reason (R) are true but Reason (R) is NOT the correct explanation of Assertion (A)
  \item (C) Assertion (A) is true but Reason (R) is false
  \item (D) Assertion (A) is false but Reason (R) is true
\end{enumerate}
\vspace{0.15cm}

\question[1] Assertion (A): The common difference of the Arithmetic Progression $5, 2, -1, -4, \dots$ is $-3$. Reason (R): The common difference $d$ of an AP is given by $d = a_{n} - a_{n-1}$.
\begin{enumerate}[label=(\Alph*)]
  \item (A) Both Assertion (A) and Reason (R) are true and Reason (R) is the correct explanation of Assertion (A)
  \item (B) Both Assertion (A) and Reason (R) are true but Reason (R) is NOT the correct explanation of Assertion (A)
  \item (C) Assertion (A) is true but Reason (R) is false
  \item (D) Assertion (A) is false but Reason (R) is true
\end{enumerate}
\vspace{0.15cm}

\question[1] Assertion (A): The common difference of the Arithmetic Progression $5, 2, -1, -4, \dots$ is $-3$.
Reason (R): The common difference $d$ of an AP is given by $d = a_{n} - a_{n-1}$.
\begin{enumerate}[label=(\Alph*)]
  \item (A) Both Assertion (A) and Reason (R) are true and Reason (R) is the correct explanation of Assertion (A)
  \item (B) Both Assertion (A) and Reason (R) are true but Reason (R) is NOT the correct explanation of Assertion (A)
  \item (C) Assertion (A) is true but Reason (R) is false
  \item (D) Assertion (A) is false but Reason (R) is true
\end{enumerate}
\vspace{0.15cm}

\question[1] Assertion (A): The common difference of the Arithmetic Progression $5, 2, -1, -4, \dots$ is $-3$. Reason (R): The common difference $d$ of an Arithmetic Progression is given by $d = a_{n} - a_{n-1}$.
\begin{enumerate}[label=(\Alph*)]
  \item (A) Both Assertion (A) and Reason (R) are true and Reason (R) is the correct explanation of Assertion (A)
  \item (B) Both Assertion (A) and Reason (R) are true but Reason (R) is NOT the correct explanation of Assertion (A)
  \item (C) Assertion (A) is true but Reason (R) is false
  \item (D) Assertion (A) is false but Reason (R) is true
\end{enumerate}
\vspace{0.15cm}

\question[1] Assertion (A): The common difference of the Arithmetic Progression $5, 2, -1, -4, ...$ is $-3$. Reason (R): The common difference $d$ of an AP is given by $d = a_{n} - a_{n-1}$.
\begin{enumerate}[label=(\Alph*)]
  \item (A) Both Assertion (A) and Reason (R) are true and Reason (R) is the correct explanation of Assertion (A)
  \item (B) Both Assertion (A) and Reason (R) are true but Reason (R) is NOT the correct explanation of Assertion (A)
  \item (C) Assertion (A) is true but Reason (R) is false
  \item (D) Assertion (A) is false but Reason (R) is true
\end{enumerate}
\vspace{0.15cm}

\question[1] Assertion (A): The common difference of the Arithmetic Progression $5, 2, -1, -4, \dots$ is $-3$. Reason (R): The common difference $d$ of an Arithmetic Progression $a_1, a_2, a_3, \dots$ is given by $d = a_n - a_{n-1}$ for $n > 1$.
\begin{enumerate}[label=(\Alph*)]
  \item (A) Both Assertion (A) and Reason (R) are true and Reason (R) is the correct explanation of Assertion (A)
  \item (B) Both Assertion (A) and Reason (R) are true but Reason (R) is NOT the correct explanation of Assertion (A)
  \item (C) Assertion (A) is true but Reason (R) is false
  \item (D) Assertion (A) is false but Reason (R) is true
\end{enumerate}
\vspace{0.15cm}

\question[1] Assertion (A): The common difference of the Arithmetic Progression 5, 2, -1, -4, ... is -3. Reason (R): The common difference $d$ of an Arithmetic Progression $a_1, a_2, a_3, ...$ is given by $d = a_{n} - a_{n-1}$.
\begin{enumerate}[label=(\Alph*)]
  \item (A) Both Assertion (A) and Reason (R) are true and Reason (R) is the correct explanation of Assertion (A)
  \item (B) Both Assertion (A) and Reason (R) are true but Reason (R) is NOT the correct explanation of Assertion (A)
  \item (C) Assertion (A) is true but Reason (R) is false
  \item (D) Assertion (A) is false but Reason (R) is true
\end{enumerate}
\vspace{0.15cm}

\question[1] Assertion (A): The common difference of the Arithmetic Progression $5, 2, -1, -4, ...$ is $-3$. Reason (R): The common difference $d$ of an AP is given by $d = a_{n} - a_{n-1}$.
\begin{enumerate}[label=(\Alph*)]
  \item (A) Both Assertion (A) and Reason (R) are true and Reason (R) is the correct explanation of Assertion (A)
  \item (B) Both Assertion (A) and Reason (R) are true but Reason (R) is NOT the correct explanation of Assertion (A)
  \item (C) Assertion (A) is true but Reason (R) is false
  \item (D) Assertion (A) is false but Reason (R) is true
\end{enumerate}
\vspace{0.15cm}

\question[1] Assertion (A): The common difference of the Arithmetic Progression $5, 2, -1, -4, ...$ is $-3$.
Reason (R): The common difference $d$ of an Arithmetic Progression is given by $d = a_{n} - a_{n-1}$.
\begin{enumerate}[label=(\Alph*)]
  \item (A) Both Assertion (A) and Reason (R) are true and Reason (R) is the correct explanation of Assertion (A)
  \item (B) Both Assertion (A) and Reason (R) are true but Reason (R) is NOT the correct explanation of Assertion (A)
  \item (C) Assertion (A) is true but Reason (R) is false
  \item (D) Assertion (A) is false but Reason (R) is true
\end{enumerate}
\vspace{0.15cm}

\question[1] Assertion (A): The common difference of the Arithmetic Progression $5, 2, -1, -4, \dots$ is $-3$. 
Reason (R): The common difference $d$ of an AP $a_1, a_2, a_3, \dots$ is given by $d = a_{n} - a_{n-1}$.
\begin{enumerate}[label=(\Alph*)]
  \item (A) Both Assertion (A) and Reason (R) are true and Reason (R) is the correct explanation of Assertion (A)
  \item (B) Both Assertion (A) and Reason (R) are true but Reason (R) is NOT the correct explanation of Assertion (A)
  \item (C) Assertion (A) is true but Reason (R) is false
  \item (D) Assertion (A) is false but Reason (R) is true
\end{enumerate}
\vspace{0.15cm}

\question[1] Assertion (A): If the $n^{th}$ term of an Arithmetic Progression is given by $a_n = 3 + 4n$, then the common difference of the AP is 4. 
Reason (R): The common difference $d$ of an AP is given by the difference between consecutive terms $a_n - a_{n-1}$.
\begin{enumerate}[label=(\Alph*)]
  \item (A) Both Assertion (A) and Reason (R) are true and Reason (R) is the correct explanation of Assertion (A)
  \item (B) Both Assertion (A) and Reason (R) are true but Reason (R) is NOT the correct explanation of Assertion (A)
  \item (C) Assertion (A) is true but Reason (R) is false
  \item (D) Assertion (A) is false but Reason (R) is true
\end{enumerate}
\vspace{0.15cm}

\question[1] Assertion (A): The sum of the first $n$ terms of an arithmetic progression is given by $S_n = 3n^2 + 5n$. The common difference of this AP is 6.
Reason (R): For an AP where $S_n = an^2 + bn$, the common difference is always $2a$.
\begin{enumerate}[label=(\Alph*)]
  \item (A) Both Assertion (A) and Reason (R) are true and Reason (R) is the correct explanation of Assertion (A)
  \item (B) Both Assertion (A) and Reason (R) are true but Reason (R) is NOT the correct explanation of Assertion (A)
  \item (C) Assertion (A) is true but Reason (R) is false
  \item (D) Assertion (A) is false but Reason (R) is true
\end{enumerate}
\vspace{0.15cm}

\question[1] Assertion (A): The common difference of the Arithmetic Progression $5, 2, -1, -4, ...$ is $-3$. 
Reason (R): The common difference $d$ of an Arithmetic Progression is calculated by the formula $d = a_{n} - a_{n-1}$ for any $n > 1$.
\begin{enumerate}[label=(\Alph*)]
  \item (A) Both Assertion (A) and Reason (R) are true and Reason (R) is the correct explanation of Assertion (A)
  \item (B) Both Assertion (A) and Reason (R) are true but Reason (R) is NOT the correct explanation of Assertion (A)
  \item (C) Assertion (A) is true but Reason (R) is false
  \item (D) Assertion (A) is false but Reason (R) is true
\end{enumerate}
\vspace{0.15cm}

\question[1] Assertion (A): The value of $\sin \theta$ is always less than 1. Reason (R): The value of $\sin \theta$ can never exceed 1 for any real value of $\theta$.
\begin{enumerate}[label=(\Alph*)]
  \item (A) Both Assertion (A) and Reason (R) are true and Reason (R) is the correct explanation of Assertion (A)
  \item (B) Both Assertion (A) and Reason (R) are true but Reason (R) is NOT the correct explanation of Assertion (A)
  \item (C) Assertion (A) is true but Reason (R) is false
  \item (D) Assertion (A) is false but Reason (R) is true
\end{enumerate}
\vspace{0.15cm}

\question[1] Assertion (A): If $\tan \theta = \cot \theta$, then $\theta = 45^\circ$. Reason (R): $\tan \theta$ and $\cot \theta$ are reciprocal to each other.
\begin{enumerate}[label=(\Alph*)]
  \item (A) Both Assertion (A) and Reason (R) are true and Reason (R) is the correct explanation of Assertion (A)
  \item (B) Both Assertion (A) and Reason (R) are true but Reason (R) is NOT the correct explanation of Assertion (A)
  \item (C) Assertion (A) is true but Reason (R) is false
  \item (D) Assertion (A) is false but Reason (R) is true
\end{enumerate}
\vspace{0.15cm}

\end{questions}
\bigskip
\noindent\textbf{\large Section B: Very Short Answer (2 marks each)}

\begin{questions}
\question[2] If $3 \cot \theta = 2$, find the value of $\frac{4 \sin \theta - 3 \cos \theta}{2 \sin \theta + 6 \cos \theta}$.
\vspace{0.15cm}

\question[2] Evaluate $\frac{\cos 45^\circ}{\sec 30^\circ + \text{cosec } 30^\circ}$.
\vspace{0.15cm}

\end{questions}
\bigskip
\noindent\textbf{\large Section C: Short Answer (3 marks each)}

\begin{questions}
\question[3] Find the sum of the first 20 terms of the AP: $1, 6, 11, 16, \dots$
\vspace{0.15cm}

\question[3] How many two-digit numbers are divisible by 3?
\vspace{0.15cm}

\question[3] If the sum of first $n$ terms of an AP is $S_n = 3n^2 + n$, find the common difference.
\vspace{0.15cm}

\question[3] The fourth term of an AP is 11 and the sum of the fifth and seventh terms is 34. Find the common difference.
\vspace{0.15cm}

\question[3] If $7 \sin^2 \theta + 3 \cos^2 \theta = 4$, show that $\tan \theta = \frac{1}{\sqrt{3}}$.
\vspace{0.15cm}

\question[3] Evaluate the following expression: $\frac{\cos 58^\circ}{\sin 32^\circ} + \frac{\sin 22^\circ}{\cos 68^\circ} - \frac{\cos 38^\circ \csc 52^\circ}{\tan 18^\circ \tan 35^\circ \tan 60^\circ \tan 72^\circ \tan 55^\circ}$.
\vspace{0.15cm}

\end{questions}
\bigskip
\noindent\textbf{\large Section D: Long Answer (4-5 marks each)}

\begin{questions}
\question[5] A sum of Rs.~1600 is to be used to give 5 cash prizes to students of a school for their overall academic performance. If each prize is Rs.~20 less than its preceding prize, find the value of each of the prizes.
\vspace{0.15cm}

\question[5] Find the sum of all three-digit numbers which leave a remainder of 3 when divided by 5.
\vspace{0.15cm}

\question[3] Find the sum of all two-digit numbers which are divisible by 6.
\vspace{0.15cm}

\question[3] If the 4th term of an A.P. is 0, prove that its 25th term is triple its 11th term.
\vspace{0.15cm}

\question[2] In an A.P., the sum of first $n$ terms is $S_n = 3n^2 + 5n$. Find the common difference of the A.P.
\vspace{0.15cm}

\question[2] Find the value of $k$ for which $k-1, k+3, 3k-1$ are in Arithmetic Progression.
\vspace{0.15cm}

\question[5] Prove the following trigonometric identity: $\frac{\cos \theta - \sin \theta + 1}{\cos \theta + \sin \theta - 1} = \csc \theta + \cot \theta$, using the identity $\csc^2 \theta = 1 + \cot^2 \theta$.
\vspace{0.15cm}

\end{questions}
\bigskip

\end{document}
